% Generated draft: requires mathematical and source review.
\documentclass[11pt]{article}
\usepackage[margin=25mm]{geometry}
\usepackage{fontspec}
\setmainfont{Latin Modern Roman}
\usepackage{amsmath,amssymb,mathtools,unicode-math}
\setmathfont{Latin Modern Math}
\usepackage{enumitem,needspace,xurl,hyperref}
\setlength{\parindent}{0pt}
\setlength{\parskip}{5pt}
\setlength{\emergencystretch}{3em}
\newcommand{\N}{\mathbb{N}}
\newcommand{\Z}{\mathbb{Z}}
\newcommand{\Q}{\mathbb{Q}}
\newcommand{\R}{\mathbb{R}}
\newcommand{\C}{\mathbb{C}}
\newcommand{\F}{\mathbb{F}}
\newcommand{\A}{\mathbb{A}}
\newcommand{\PP}{\mathbb{P}}
\newcommand{\E}{\mathbb{E}}
\newcommand{\eps}{\varepsilon}
\newcommand{\dd}{\,\mathrm{d}}
\begin{document}
\section*{UnsolvedMath ID: 30005124}
\textbf{Source problem code:} catalogue entry
\subsection{Definitions and mathematical statement}
The supplied record concerns cohomological Hall algebras, denoted by \(A_{Q,W}\), associated with a quiver \(Q\) and a potential \(W\). It does not specify the construction or conventions for these algebras. Its recoverable question is the following schematic existence problem: for the intended reference pair \((\widetilde Q,\widetilde W)\), does there exist a quiver \(Q'\) with potential \(W'\) such that
\[
A_{Q',W'}\text{ ``partially fermionizes'' }A_{\widetilde Q,\widetilde W}?
\]
The reference pair and the phrase ``partially fermionizes'' are not defined in the supplied record. In particular, no required relation between the two algebras, their gradings, or their BPS cohomology is specified. Consequently, this expression cannot yet be made into a precise mathematical predicate.

The extracted background supplies the following intended motivation, with essential geometric data missing. For every \(\alpha\in\mathfrak h\), including \(\alpha=0\), it refers to a Cartesian construction of a smooth threefold \(Y^\alpha\). The diagram defining this construction is corrupted. For a vertex \(i\in\Gamma_0\), it states that a corresponding curve \(C_i\) deforms along \(\mathbb A^1\) if and only if \(\alpha\in i^\perp\). Here \(\Gamma\), \(\mathfrak h\), and the hyperplanes \(i^\perp\) are not defined in the excerpt. Writing \(M_{i,n}\) for the moduli space of semistable coherent sheaves of the class denoted by \(\mathcal O_{C_i}(n)\), the source reports
\[
M_{i,n}\cong
\begin{cases}
\mathbb A^1,&\alpha\in i^\perp,\\
\mathrm{pt},&\alpha\notin i^\perp.
\end{cases}
\]
The range of \(n\), the meaning of the class, the stability condition, and the ambient moduli problem are not supplied. The source describes these moduli spaces as smooth and simply connected and writes their BPS cohomology as
\[
\mathfrak g^{\mathrm{BPS}}_{i,n}
=H\!\left(M_{i,n},\mathbb Q[\dim(M_{i,n})]\right)[-1].
\]
It reports that this cohomology is one-dimensional, with parity depending on which hyperplanes \(\alpha\) avoids. It also cautions that BPS cohomology for a specified class in a specified \(3\)-Calabi--Yau category can in general be difficult to calculate or even define. The displayed formula is therefore retained as source-specific background, without extending it to arbitrary such categories.

For its connection with the noncommutative conifold, the source specializes to \(G=\mathbb Z/2\mathbb Z\), states \(\mathfrak g=\mathfrak{sl}_2\) and \(\dim\mathfrak h=1\), and refers to the extended quiver associated with \(\Gamma\). It reports a derived equivalence between \(\operatorname{Coh}(Y^0)\) and the category of \(\Pi_Q[\omega]\)-modules, where \(Y^0\) is the specialization at \(\alpha=0\). The intended derived categories and the algebra \(\Pi_Q[\omega]\) are not defined. It further states that the BPS cohomology for \((\widetilde Q,\widetilde W)\) is entirely even. The sentence concerning every \(\alpha\ne0\) and the resolved conifold is truncated, so its full conclusion cannot be recovered.

The supplied registry marks this statement as \texttt{unrecoverable} and records it as a duplicate extraction of \texttt{OWR-10252931-001}, with no separate clean statement promoted. The formulation above preserves the recoverable question and background; it is a statement-only draft requiring human mathematical review.
\subsection{Short English statement}
The recoverable question asks whether a quiver with potential can produce a cohomological Hall algebra that ``partially fermionizes'' the intended reference algebra. The meaning of this operation and the reference pair are missing. The registry identifies this corrupted, truncated record as a duplicate of \texttt{OWR-10252931-001}, so a precise independent statement cannot be recovered from the supplied evidence.
\subsection{Sources}
\begin{itemize}
\item Local UnsolvedMath registry, permanent ID \(30005124\), title: \emph{Partial Fermionization of Cohomological Hall Algebras}; exact source problem code: \texttt{catalogue entry}. The record belongs to the Oberwolfach Reports Open Problems collection and gives proposed year \(2022\); no individual proposer is supplied.
\item Supplied report citation: \emph{Interactions between Algebraic Geometry and Noncommutative Algebra (2023)}, Oberwolfach Reports, pp.~1239--1312, DOI \texttt{10.4171/owr/2022/23}. Supplied source URL and literature-source link: \url{https://doi.org/10.4171/owr/2022/23}. The supplied background describes these reports as documenting mathematical workshops and their open-problem discussions. This citation has not been independently checked here.
\item Supplied literature source: ``Davison, fermionization paper,'' \url{https://doi.org/10.1017/S001708952200009X}. No fuller bibliographic details are supplied.
\item Supplied literature assessment dated exactly \texttt{2026-08-22}: the status is labeled ``partially solved.'' The assessment says that a concrete fermionization procedure is known for important preprojective-algebra CoHAs, including central-extension and parity results, but that no general quiver-and-potential construction answering the question is known. This is supplied dated curation derived from the OWR record, not an independently verified mathematical or literature assessment. The supplied background explicitly requests independent verification of time-sensitive status and cited claims.
\item Supplied statement-verification note: the official report is said to confirm that this is a duplicate of the quiver-with-potential formulation retained as \texttt{OWR-10252931-001}; no separate clean statement is promoted, and the original remains visible for provenance. This verification claim is part of the supplied registry data and has not been independently checked here.
\end{itemize}
\end{document}
